\section*{Chapter \ref{ch:b3:regression-discontinuity} --- Regression Discontinuity}

\begin{solution}{exo:b3:regression-discontinuity:1}
The left line reaches 0.4597 at the \omterm{def:b3:regression-discontinuity:running-variable}{cutoff} and the right line 0.5353; the jump is
$0.5353 - 0.4597 = 0.0756$. % figdata: elections_fit_left, elections_fit_right
Margins are drawn mostly near zero, so the bins near $\pm 0.4$ hold few
elections and their means are noisier.
\end{solution}

\begin{solution}{exo:b3:regression-discontinuity:2}
40: one class of 40. 41: $41/2 = 20.5$. 80: $80/2 = 40$. 81: $81/3 = 27$. 121:
$121/4 = 30.25$. The rule jumps at 41 (from 40 to 20.5, by $-19.5$), at 81 (from 40
to 27, by $-13$) and at 121 (from $120/3 = 40$ to 30.25, by $-9.75$); the jumps
shrink because the extra class is shared among more pupils.
\end{solution}

\begin{solution}{exo:b3:regression-discontinuity:3}
$52 - 40 = 12$ percentage points. It assumes that applicants on 68--69 are
comparable to those on 70--71 apart from the scheme. If employment rises with the
score, the two groups differ by about two marks of that slope as well: the
difference of means adds the slope times two marks to the effect.
\end{solution}

\begin{solution}{exo:b3:regression-discontinuity:4}
That pupils on the two sides differ in background, so something other than the
scholarship changes at the \omterm{def:b3:regression-discontinuity:running-variable}{cutoff}: sorting (better-connected families obtain
re-marks), or another programme using the same mark. Look at the density of
scores at the \omterm{def:b3:regression-discontinuity:running-variable}{cutoff}, at other predetermined covariates, at the rules for
re-marking, and at whether any other benefit uses the same mark.
\end{solution}

\begin{solution}{exo:b3:regression-discontinuity:5}
The tested solution is \code{ex05_global_polynomial.py}, in the chapter's
\code{exercises/} folder. Bias and standard deviation: order 1, $-0.034$ and
0.010; order 2, 0.010 and 0.012; order 4, $-0.001$ and 0.022; order 6, $-0.002$
and 0.029; local linear within 0.25, $-0.003$ and 0.013. % output: ex05_global_polynomial
Low orders are biased because they cannot follow the curve; high orders are
unbiased here but noisy, because observations far from the \omterm{def:b3:regression-discontinuity:running-variable}{cutoff} pull the
fitted value at the \omterm{def:b3:regression-discontinuity:running-variable}{cutoff} about. The local linear fit has a small bias and
half the noise of the order-6 polynomial. Here the true curve is a cubic, which
an order-4 fit contains; with real data no order is known to be right.
\end{solution}

\begin{solution}{exo:b3:regression-discontinuity:6}
\begin{enumerate}
\item $(0.75 - 0.70)/(0.85 - 0.60) = 0.05/0.25 = 0.20$: insurance raises the
probability of a visit by 20 percentage points, for the people at the threshold
who are insured because they crossed it.
\item Those insured below the threshold are \omterm{def:b3:instrumental-variables:compliance}{always-takers}: crossing does not
change their insurance, so they contribute nothing to either jump, and the ratio
says nothing about them.
\end{enumerate}
\end{solution}

\begin{solution}{exo:b3:regression-discontinuity:7}
Upward. The determined pupils move from just below the \omterm{def:b3:regression-discontinuity:running-variable}{cutoff} to just above it,
raising later earnings above the \omterm{def:b3:regression-discontinuity:running-variable}{cutoff} and lowering them below it for reasons
other than the scholarship.
\end{solution}

\begin{solution}{exo:b3:regression-discontinuity:8}
The tested solution is \code{ex08_donut.py}, in the chapter's \code{exercises/}
folder. With all recorded scores the jump is 0.076 (s.e.\ 0.015); dropping the
marks 69 and 70 gives 0.089 (0.022); dropping 68--71 gives 0.087 (0.028).
% output: ex08_donut
The re-marked pupils, all recorded at 70, have the earnings of pupils on 66--69,
so they pull the mean just above the \omterm{def:b3:regression-discontinuity:running-variable}{cutoff} down and the estimate below the true
0.10. The donut removes them; the estimate moves toward the truth, and its
standard error rises because the observations closest to the \omterm{def:b3:regression-discontinuity:running-variable}{cutoff}, the most
informative, are gone.
\end{solution}

\begin{solution}{exo:b3:regression-discontinuity:9}
A heap of households just below the line and a gap just above it, or households
just below it that differ in their characteristics (connections to the council,
a particular district) from those just above. If the line was chosen to include
particular households, the units on the two sides are not alike, continuity
fails, and the jump mixes the subsidy with the reasons those households were
chosen.
\end{solution}

\begin{solution}{exo:b3:regression-discontinuity:10}
The OLS intercept of $Y$ on $(1, x)$ is $\hat a_1 = (S_2 \sum_i Y_i - S_1 \sum_i
x_i Y_i)/(n_1 S_2 - S_1^2)$, from the normal equations, which is $\sum_i
\omega_i Y_i$ with $\omega_i = (S_2 - S_1 x_i)/(n_1 S_2 - S_1^2)$. The weights sum
to $(n_1 S_2 - S_1^2)/(n_1 S_2 - S_1^2) = 1$. The denominator is positive (it is
$n_1$ times the sum of squared deviations of $x$). Above the \omterm{def:b3:regression-discontinuity:running-variable}{cutoff} all $x_i \ge
0$, so $S_1 > 0$, and $\omega_i < 0$ exactly when $x_i > S_2/S_1$: the
observations farthest from the \omterm{def:b3:regression-discontinuity:running-variable}{cutoff} get negative weight, as the line is
extrapolated back to $x = 0$.
\end{solution}

\begin{solution}{exo:b3:regression-discontinuity:11}
The tested solution is \code{ex11_coverage.py}, in the chapter's
\code{exercises/} folder. Coverage is 0.944 at a window of 0.02, 0.932 at 0.10,
0.882 at 0.25 and 0.018 at 0.50. % output: ex11_coverage
The interval is built for noise only. At 0.50 the bias is $-0.015$ and the
standard deviation 0.004 (\cref{out:b3:regression-discontinuity:bandwidth}): the
interval is centred nearly four standard errors from the truth and almost never
covers it. At 0.25, where the error-minimising window lies, the bias is already a
large part of the standard error and coverage falls short of 95~\%.
\end{solution}

\begin{solution}{exo:b3:regression-discontinuity:12}
The pupils at 50 differ from those at 70: they may gain more from funding (if
money, not ability, is what holds them back) or less (if they are less
prepared for the course the scholarship pays for). And a scholarship given to
many more pupils may change what it pays for: crowded courses, a weaker signal
to employers. A design: if the \omterm{def:b3:regression-discontinuity:running-variable}{cutoff} differed between years or towns, a
discontinuity at 50 where the \omterm{def:b3:regression-discontinuity:running-variable}{cutoff} was 50 estimates the effect there; or a
lottery among applicants between 50 and 70.
\end{solution}

\begin{solution}{pb:b3:regression-discontinuity:scholarship}
\textbf{1.} Sharp for the offer, which goes to everyone from 70 and to no one
below; fuzzy for the award, which jumps by about two thirds.
% figdata: scholarship

\textbf{2.} Being offered the scholarship raises log earnings at 30 by 0.0792,
about 8~\%, for pupils at the \omterm{def:b3:regression-discontinuity:running-variable}{cutoff}.

\textbf{3.} $0.0792 \pm 1.96 \times 0.0323 = 0.0792 \pm 0.0633$: from 0.0159 to
0.1425. It excludes zero.

\textbf{4.} 0.0784, 0.0792 and 0.0738: the estimate barely moves, so curvature
is not driving it. The standard errors, 0.0442, 0.0323 and 0.0241, fall because a
wider window uses more pupils; doubling the window divides them by about
$\sqrt 2$ ($0.0442/0.0323 = 1.37$).

\textbf{5.} $-0.0101$ with a standard error of 0.0309: no jump in a variable
fixed before the exam. Pupils just above and just below are alike in family
background, as the design requires.

\textbf{6.} Only part of those offered take an award because of the offer; the
effect of the offer is the effect of the award diluted by that share.

\textbf{7.} Crossing 70 raises the probability of receiving an award by 0.6675.

\textbf{8.} $0.0792/0.6675 = 0.119$; $\eu^{0.119} - 1 = 0.126$: earnings 12.6~\%
higher.

\textbf{9.} $0.0784/0.6552 = 0.120$ and $0.0738/0.6525 = 0.113$.

\textbf{10.} $0.0323/0.6675 = 0.0484$; $0.119 \pm 1.96 \times 0.0484 = 0.119 \pm
0.0949$, from 0.024 to 0.214. The approximation ignores the noise in the
award jump, which is small here.

\textbf{11.} For \omterm{def:b3:instrumental-variables:compliance}{compliers} at 70: pupils who receive an award if and only if
they are offered one. It leaves out \omterm{def:b3:instrumental-variables:compliance}{never-takers}, who decline even when
offered, and \omterm{def:b3:instrumental-variables:compliance}{always-takers}, who obtain an equivalent award below 70.

\textbf{12.} Crossing 70 changes earnings only by changing whether a pupil
receives the award. It fails if the offer itself matters: if pupils list it on
their applications and employers reward it, or if being offered raises a
pupil's ambition even when the award is declined.

\textbf{13.} 0.0211 (0.0277), $t = 0.76$: no jump where nothing changes, so the
local linear fits are not manufacturing jumps out of the curve's shape.

\textbf{14.} A window that crossed 70 would include the real jump in the fit on
one side and could turn it into a spurious jump at the placebo \omterm{def:b3:regression-discontinuity:running-variable}{cutoff}.

\textbf{15.} At 10, with the estimates at 5 and 20 beside it, the balance test,
the placebo \omterm{def:b3:regression-discontinuity:running-variable}{cutoff} and the density test --- and the fuzzy estimate with its
interval, since the award, not the offer, is the policy's instrument.

\textbf{16.} A cohort or a town where the mark of 70 carried no scholarship (or
the scholarship's \omterm{def:b3:regression-discontinuity:running-variable}{cutoff} was elsewhere) but the same standing with employers: a
jump at 70 there would be the signal alone.

\textbf{17.} All pupils: $0.3036/0.0798 = 3.80$. Without the school: $0.0623/0.0841
= 0.74$. % figdata: scholarship_density
Too many pupils just above 70, all coming from the one school; elsewhere the
density is smooth.

\textbf{18.} $\eu^{0.3036} = 1.35$: about 35~\% more pupils per mark just above
70 than just below.

\textbf{19.} Pupils of that school whose true score was 67, 68 or 69. They are
on average about two marks weaker than the honest pupils on 70.

\textbf{20.} Downward, since weaker pupils are added just above the \omterm{def:b3:regression-discontinuity:running-variable}{cutoff}; but
the bias should be small, as few pupils were moved and they were only two marks
weaker. The estimates without the school are 0.0628, 0.0770 and 0.0729, against
0.0784, 0.0792 and 0.0738: they are not higher, and the differences are well
inside a standard error. The re-marking is detected but does no visible harm.

\textbf{21.} Parents' income moves little with the score, and the moved pupils
are few and only slightly weaker: the jump they cause in mean parents' income is
far below what a standard error of 0.03 can detect.

\textbf{22.} If that school's pupils differ from the rest (a richer district), the
population at the \omterm{def:b3:regression-discontinuity:running-variable}{cutoff} changes, and the estimate is for a different group; and
if other teachers re-marked without being caught, dropping one school does not
remove the problem.

\textbf{23.} Drop the pupils recorded on 69 and 70 (or on 68 to 71) and estimate
from the rest; the re-marked pupils, all on 70, are gone, at the cost of the
observations nearest the \omterm{def:b3:regression-discontinuity:running-variable}{cutoff}.

\textbf{24.} The scholarship raises earnings at 30 by 0.119 log points, 12.6~\%,
for pupils at the \omterm{def:b3:regression-discontinuity:running-variable}{cutoff} of 70 whose award depends on the offer, with an
approximate 95~\% interval of 0.024 to 0.214. The density test finds re-marking
in one school ($t = 3.80$) and none elsewhere ($t = 0.74$); without that school
the estimate is $0.0770/0.6613 = 0.116$.
\end{solution}
